% ICASSP 2027 submission -- DRAFT
% Deadline: 2026-09-16 (UTC-12). 4 pages + 1 page references.
% Requires the official ICASSP template (spconf.sty, IEEEbib.bst).
%
% STATUS: every number marked [V] below is verified in-repo; see
% paper/icassp2027/NUMBERS.md for the provenance of each one.
\documentclass{article}
\usepackage{spconf,amsmath,graphicx,booktabs,microtype}

\title{The Synthetic-to-Real Gap in Sheet-Image Score Following\\
       is Representational, not Acoustic}

\name{Anonymous}
\address{Anonymous}

\begin{document}
\maketitle

\begin{abstract}
Score followers that track a live performance on a sheet-music image degrade
sharply when moved from synthesised audio to real recordings. This degradation
is routinely attributed to acoustic mismatch --- reverberation, microphone
response, instrument timbre --- and is addressed with acoustic remedies. We
show that this attribution is largely wrong. Using two microphones placed on
the \emph{same} performances, so that the performance is held exactly fixed, we
measure that room reverberation costs a stock automatic music transcription
(AMT) system essentially nothing: onset $F_1$ is $0.9116$ on the room
microphone versus $0.9124$ on a direct pickup at $50\,$ms tolerance. On those
same recordings, image-conditioned trackers lose roughly thirty points of
tracking accuracy. The note-level information survives the room almost intact;
the tracking architecture discards it. We corroborate this with an intervention
asymmetry: augmentation with real measured impulse responses, drawn from one
bank and applied under one protocol, is reported to be worth $+25.2$ points to
a detection-based tracker but is worth only $+11.0$ to a dense-heatmap tracker.
We argue the binding constraint is the output parameterisation rather than the
audio front end, and we report three measurement problems --- an aggregation
mismatch worth $+7.3$ points, an effective sample size of $25$ rather than
$4415$, and undisclosed oracle inputs --- that make small reported gains on
this benchmark uninterpretable.
\end{abstract}

\begin{keywords}
score following, music alignment, optical music recognition, evaluation,
domain shift
\end{keywords}

\section{Introduction}
\label{sec:intro}

Sheet-image score following takes a live audio performance and a rendered image
of the score, and predicts the performer's current position on the image, frame
by frame. Unlike symbolic score following, it consumes the score as pixels,
which removes any need for a machine-readable encoding of the piece.

Systems in this family are trained almost exclusively on synthesised audio,
because ground truth requires an exact alignment between audio and notehead
coordinates, and that is cheap to obtain only by rendering. Every such system
degrades when evaluated on real recordings. The degradation is large: on the
real-audio subset of MSMD, published CUNet scores $22.4$ where the same
architecture reaches the high eighties on synthetic audio.

The prevailing explanation is acoustic. Real rooms add reverberation, real
microphones colour the spectrum, and real pianos differ from the soundfont, so
the audio encoder sees inputs unlike anything in training. This explanation
motivates acoustic remedies --- impulse-response augmentation, multi-condition
training, test-time normalisation --- and it is the framing under which we
ourselves worked for several months.

This paper argues that the explanation is mostly incorrect, and that acting on
it has misdirected effort.

Our contributions:
\begin{enumerate}
\itemsep0em
\item A controlled measurement showing that the room preserves note-level
      information almost perfectly (Section~\ref{sec:twomic}). Two microphones
      on the same take isolate the room with the performance held fixed.
\item An intervention asymmetry showing that the same acoustic remedy pays very
      differently depending on the output parameterisation
      (Section~\ref{sec:where}).
\item Three measurement problems on this benchmark that jointly make gains
      below roughly ten points uninterpretable (Section~\ref{sec:eval}), with
      a corrected reporting protocol.
\end{enumerate}

\section{Setup}
\label{sec:setup}

\noindent\textbf{Data.} We use MSMD and its real-recording extension. We
evaluate on four acoustic tiers over the same $25$ pages:
\emph{synth} (soundfont rendering, matching training);
\emph{rp\_synth} (a matched synthetic control: same pages, same ground truth,
training soundfont);
\emph{di-left} (real piano, direct pickup, no room);
and \emph{room} (real piano, room microphone). Because \emph{di-left} and
\emph{room} are two microphones recording the \emph{same} take, their
difference isolates the room with the performance held exactly fixed. This
pairing is the instrument used throughout Section~\ref{sec:twomic}.

\noindent\textbf{Metric.} We report \textsc{pct@0.5s}: the percentage of onsets
whose predicted position, mapped back through the score-to-onset interpolation,
falls within $0.5\,$s of the true onset. Thresholds are
$\{0.05, 0.1, 0.5, 1.0, 5.0\}\,$s. All numbers are onset-only and
\emph{micro}-aggregated (onsets pooled across pieces, then divided), matching
the convention of the baselines we compare against. Section~\ref{sec:eval}
shows why that choice must be stated.

\noindent\textbf{Baselines.} We reproduced the strongest reproducible published
system ourselves rather than quoting it, obtaining
$63.0$ at $0.1\,$s and $79.9$ at $0.5\,$s on \emph{room}, which matches the
published row exactly. Numbers frequently quoted as $63.0$ for this system are
the $0.1\,$s column; the $0.5\,$s column is $79.9$.

\section{Does the room destroy the information?}
\label{sec:twomic}

If the synthetic-to-real gap were acoustic in the strong sense --- the room
degrading the audio so that note events are no longer recoverable --- then any
system reading notes from that audio should degrade too. We test this directly
with an off-the-shelf AMT system, which shares no components with our tracker
and was trained on entirely different data.

We transcribe all $25$ real performances in both microphone conditions with two
AMT models: a widely used piano transcription model (\textsc{kong}) and a
reverb-augmented variant (\textsc{edwards}). We score note onsets with
\texttt{mir\_eval} at $50$ and $100\,$ms.

\begin{table}[t]
\centering
\caption{Note-onset $F_1$ on the same $25$ performances, two microphones,
         same take. The room costs a transcriber essentially nothing.}
\label{tab:twomic}
\begin{tabular}{llccc}
\toprule
model & tol & room & di-left & $\Delta$ \\
\midrule
\textsc{kong}     & $50\,$ms  & 0.9116 & 0.9124 & $-0.001$ \\
\textsc{kong}     & $100\,$ms & 0.9546 & 0.9903 & $-0.036$ \\
\textsc{edwards}  & $50\,$ms  & 0.9441 & 0.9259 & $+0.018$ \\
\textsc{edwards}  & $100\,$ms & 0.9864 & 0.9909 & $-0.005$ \\
\bottomrule
\end{tabular}
\end{table}

Table~\ref{tab:twomic} is the central measurement of this paper. At $50\,$ms
the room costs \textsc{kong} $0.001$ $F_1$, and the reverb-augmented model is
in fact \emph{better} on the room microphone than on the direct pickup. At
$100\,$ms the largest penalty is $0.036$. On these same files, our trackers
lose roughly thirty points of \textsc{pct@0.5s}.

\noindent\textbf{Validating the measurement path.} Before trusting these
numbers we verified that the coordinate round-trip is lossless. Feeding the
ground-truth MIDI in as if it were a transcription, and decoding position
through the same score-note-to-pixel-to-onset-frame path used for the AMT
output, yields exactly $100.00$ \textsc{pct@0.5s} when the tracker is allowed
to hear the note it is scored on, and $94.79$ under the stricter convention
where it is not. A coordinate or indexing error anywhere in this path would
show up as a ceiling below $100$.

\noindent\textbf{What this does and does not show.} It shows the note
information survives the room. It does \emph{not} show that a decomposed
transcribe-then-align system is a better score follower. Offline dynamic time
warping over these transcriptions reaches $98.06$ on \emph{room}, but that
figure is not comparable to an online tracker: it is non-causal, and it is
close to tautological on this data, because the recordings are reproducing-piano
playback of the score MIDI and therefore contain no tempo deviation for the
alignment to resolve. The corresponding \emph{online} greedy matcher reaches
only $10.7$, with a median error of eight to thirteen seconds; it loses the
pointer and never recovers. Causal matching, not transcription, is the hard
part of that pipeline.

\section{Where the gap actually is}
\label{sec:where}

\subsection{The same remedy pays differently}

If the audio is not the problem, acoustic augmentation should help less than
expected --- and should help \emph{unevenly} across architectures. It does.

Augmenting with real measured impulse responses is reported to move a
detection-based tracker from $46.0$ to $71.2$ on \emph{room}, a gain of
$+25.2$. We applied augmentation from a bank of $693$ real measured impulse
responses to our own dense-heatmap tracker, re-encoding degraded waveforms
through the frozen audio front end so the degradation cannot be undone by a
normalisation layer. The gain was $+11.0$ ($45.6 \rightarrow 56.6$): real, the
largest single improvement we have obtained, and well under half of what the
same class of intervention buys a detection model.

\noindent\textbf{A negative result we must disclose.} Our own earlier
conclusion that impulse-response augmentation does not help was an artefact. The
convolution used \texttt{mode='same'}, which advances the signal by
$(\lvert h \rvert - 1)/2$ samples relative to its onset labels --- a
$4$ to $20$ frame desynchronisation on half of all augmented samples. The
experiment that produced the $+11.0$ above differs from the discarded one
chiefly in that this bug is fixed and the impulse responses are measured rather
than synthesised. We report this because the failure mode is silent: the
training loss is unremarkable and only the real-audio metric suffers.

\subsection{The output parameterisation}

A controlled comparison in the literature points to the same place. Within one
paper, on one dataset, with the same audio front end, a system that predicts
position as a \emph{bucketed softmax classification} scores $58.5$ on
\emph{room} where the dense soft-Dice heatmap formulation scores $22.4$ --- a
$36$-point difference attributable to the output layer alone. This is
consistent with the segmentation literature, where soft-Dice is known to
perform well in-distribution and poorly out-of-distribution, and
synthetic-to-real is exactly such a shift.

Three independent observations therefore converge: the audio retains the
information (Section~\ref{sec:twomic}); acoustic remedies transfer only partly
to heatmap models; and the largest controlled swing in the literature is
attributable to the output parameterisation. We take this as the strongest
available guidance for where effort in this task should go.

\section{Measurement problems on this benchmark}
\label{sec:eval}

Alongside the above we found three properties of the standard evaluation that
affect how any result on it should be read.

\noindent\textbf{Aggregation.} Reported numbers in this literature use two
different estimators without stating which. Pooling all onsets and dividing
(\emph{micro}) and averaging per-piece percentages (\emph{macro}) are not
interchangeable: on one unchanged checkpoint we measure $45.4$ micro, $46.8$
macro over pages, and $52.7$ macro over pieces. The spread of $+7.3$ points
arises because a single difficult nocturne constitutes $29.8\%$ of the micro
metric but only $6.2\%$ of the piece-macro metric, and it is the recording on
which every system we have tested performs worst. Macro aggregation
systematically discounts the hardest part of the benchmark.

\noindent\textbf{Effective sample size.} Onsets within a piece are not
independent. A piece-level cluster bootstrap gives a design effect of
approximately $180$, so the effective sample size of the real-audio benchmark
is about $25$ pieces, not $4415$ onsets. The minimum detectable difference is
roughly ten points. We confirmed the practical consequence directly: retraining
one recipe twice, changing nothing but the random seed and the run, produced
$22.7$ and $16.4$ --- a spread larger than many published ablation effects.

\noindent\textbf{Undisclosed oracle inputs.} Systems differ in what structure
they are given versus what they must infer. In one recent system, ground-truth
system and bar bounding boxes are required arguments of the forward pass at
both training and inference; its loss contains no localisation term of any
kind. It selects among given regions, whereas the baseline it is compared
against must produce those regions from the image. This is not stated in the
paper and is visible only in the released code.

\noindent\textbf{Recommended protocol.} We suggest that work on this benchmark
report (i) both micro and macro aggregation, (ii) piece-level bootstrap
confidence intervals rather than point estimates, (iii) results across acoustic
tiers rather than a single condition, including a matched synthetic control,
and (iv) an explicit statement of which structures the model receives as input
and which it must predict.

\section{Conclusion}
\label{sec:conclusion}

The synthetic-to-real gap in sheet-image score following is not primarily an
acoustic problem. A transcriber recovers $91$--$95\%$ of onsets from the same
room recordings on which trackers lose thirty points, and loses essentially
nothing to the room when the performance is held fixed. The information is
present in the audio and is discarded downstream. Consistent with this,
acoustic augmentation transfers only partially to dense-heatmap models, and the
largest controlled effect in the literature is attributable to the output
parameterisation. We also find that the benchmark, as commonly used, cannot
support the size of the improvements typically claimed on it.

\vfill\pagebreak

\bibliographystyle{IEEEbib}
\bibliography{refs}

\end{document}
